\chapter{Nội Quy Và Giấy Kiểm Điểm}
\chapterintro{Giấy kiểm điểm là thể loại văn học mà không học sinh nào muốn luyện nhiều}{Nhưng vì nó quá quen thuộc nên không thể vắng trong sách này.}

\begin{lucbatbox}{Lục bát: Viết bản kiểm điểm}
Hôm nay em viết rất lâu
Từng câu xin lỗi từng đầu nhận sai
Mực còn chưa kịp khô ngay
Bạn bè đã hỏi lần này lỗi chi
Em rằng lỗi cũng thường thôi
Chỉ vì vui quá quên lời dặn cô
\end{lucbatbox}

\begin{tutuyetbox}{Thất ngôn tứ tuyệt: Nội quy rất hiền}
Nội quy treo đó hiền lành thôi
Đọc qua tưởng chuyện ở xa rồi
Đến khi ngồi viết dòng phân giải
Mới hay từng chữ dính thân mình
\end{tutuyetbox}

\begin{ghichu}
Giấy kiểm điểm thường không làm học sinh sợ bằng việc phải đọc lại chính lỗi của mình bằng chữ viết tay thật rõ.
\end{ghichu}

\ngancach

\begin{lucbatbox}{Lục bát: Hứa sửa rất chân thành}
Cuối bản em viết một câu
Từ nay xin hứa nguyện mau đổi mình
Lời văn tha thiết chân tình
Đọc lên như thể học sinh hóa hiền
Chỉ mong sau bản kiểm liền
Em nhớ lời hứa lâu hơn một tuần
\end{lucbatbox}

\begin{tutuyetbox}{Thất ngôn tứ tuyệt: Ký tên vào lỗi}
Ký tên xong thấy nặng trong đầu
Giấy trắng mà ghì cả nỗi sầu
Nếu vui bớt quá vài ba bận
Đâu phải ngồi ghi lỗi bấy lâu
\end{tutuyetbox}